\section{Scalar concentration and polynomial fitting}\label{sec:scalar-tools}

We prove two consequences of scalar concentration.  First, a small
spacetime support forces a heavy bin of quadratic trajectory
coefficients.  Second, when a common row subtracts the tangential
coordinates from a quadratic signal, concentration of the residual
scalar forces polynomial fits to that row along labeled trajectories.
The second argument begins with entropy estimates under simultaneous
time tests.  A planar dot-product theorem then gives affine row fits
and first-jet information; divided differences raise the jet order,
and a separate time test produces the polynomial fits.

The arguments in this section take place in tangent units.
Thus \(K=N^{o(1)}\) denotes a changeable subpower factor, and a measure is
\emph{dense} if its mass is at least \(K^{-1}\).  All coefficient ranges
and all fixed-dimensional norms below are bounded by \(K\).
The time interval is \([0,1]\).  A trajectory index may include auxiliary
data; a bin of its coefficients never bins those additional indices.

We use measurable time labels \(E_l\subset[0,1]\).
An instantaneous labeled mass is
\[
 \nu\{l:t\in E_l,\ \text{the indicated conditions hold}\};
\]
it is not divided by the success probability at time \(t\).
This convention also covers an incidence density \(w(l,t)\le K_{\rm pre}\)
with \(K_{\rm pre}=N^{o(1)}\), relative to \(\nu\otimes dt\).  Indeed, if its mass is
\(\delta\ge K_{\rm pre}^{-1}\), deleting \(w<\delta/2\) loses at most
\(\delta/2\).  On the remaining indicator label set, the two measures
compare by subpower factors, and its product measure is at least
\(\delta/(2K_{\rm pre})\).  The same argument applies in typical conditional
worlds.  Restrictions made in this section change the labels and leave
the indicated trajectory prior fixed unless a new normalization is
explicitly stated.

\subsection{Concentration of quadratic coefficients}

\begin{lemma}[Scalar coefficient clustering]\label{lem:scalar-clustering}
Let
\[
 x_l(t)=x_{0,l}+t x_{1,l}+t^2x_{2,l},
\]
and let the trajectory prior \(\nu\) have total mass at most \(K\).
Fix \(0\le d<1\) and a resolution \(p=N^{-c+o(1)}\), where \(c>0\)
is fixed.  Suppose that a dense labeled incidence measure, dominated
by \(K\,d\nu\,dt\), occupies at most \(Kp^{-1-d}\) squares of side \(p\)
in the \((t,x)\)-plane.  Suppose also that
\begin{equation}\label{a04:coefficient-cap}
 \nu\{l:(x_{0,l},x_{1,l},x_{2,l})\in B\}\le K r^d
\end{equation}
for every coefficient ball \(B\) of radius \(r\ge p\) in the bounded
coefficient range.  Then some coefficient bin of side \(p\) has
trajectory mass at least \(p^d/K\).
The conclusion can be obtained using only trajectories with labeled
length at least \(K^{-1}\), and it applies anew to every dense residual
satisfying the same upper bounds.
\end{lemma}

\begin{proof}
First discard trajectories of labeled length below a sufficiently small
reciprocal subpower.  The remaining trajectory mass is at least \(K^{-1}\),
and each retained trajectory meets at least \(p^{-1}/K\) labeled time
bins.  Adjoin an independent uniform variable \(D\in[0,1]\) and lift the
signal to the affine planar motion
\begin{equation}\label{a04:quadratic-lift}
 z_{l,D}(t)=\bigl(x_{0,l}+t(x_{1,l}-D),\,D+t x_{2,l}\bigr).
\end{equation}
The identity
\[
 (1,t)\cdot z_{l,D}(t)=x_l(t)
\]
shows that the marked lifted motions occupy at most \(Kp^{-2-d}\)
spacetime cubes.  In fact, above one occupied \((t,x)\)-square there
are at most \(K/p\) second-coordinate bins, and then \(K\) possible
first-coordinate bins.  Motion within one time bin changes this count
only by a subpower factor.

Apply \Cref{lem:weighted-planks}, namely the weighted full-trace form of
\citet[Lemma~2.3]{ThreeDimensional}.  Rescaling a subpower bounded spatial
chart into a fixed bounded chart and rounding line parameters much more
finely than \(p\) are allowed by that lemma.  The total weighted number
of marks is at least \(p^{-1}/K\), whereas the occupied-cell count is at
most \(Kp^{-2-d}\).  Their ratio is at least \(p^{1+d}/K\).
Consequently there is a full-time plank with fixed orthogonal spatial
axes, affine center, and physical half-widths
\[
 p/K\le u\le v\le K
\]
whose lifted trajectory mass is at least
\begin{equation}\label{a04:plank-lower}
 K^{-1}p^{d-1}uv.
\end{equation}
Here the conversion from widths in cell units contributes \(p^{-2}\);
the fixed positive loss in the quoted estimate is sent to zero slowly
after its scale is fixed.

All scalar quadratics arising from this plank lie in one coefficient
ball of radius \(Kv\).  To see this, take the difference of two lifted
motions in the plank.  Its scalar projection by \((1,t)\) has size
\(O(v)\) throughout unit time, and evaluation at three fixed separated
times bounds its three coefficients by \(O(v)\).

For fixed scalar coefficients, varying \(D\) changes the lifted motion
by \(D(-t,1)\).  If \(n=(n_1,n_2)\) is the unit narrow normal, then
\[
 \max\{|n_2|,|n_2-n_1|\}\ge c
\]
for an absolute \(c>0\).  The narrow constraint at one of \(t=0,1\)
therefore accepts a set of \(D\)'s of length at most \(Cu\).
The coefficient cap implies that the lifted plank mass is at most
\(Kuv^d\).  When \(v<p\), use the cap at \(p\) and
\(v\ge p/K\), absorbing the resulting subpower factor into \(K\).
Comparing with \Cref{a04:plank-lower} gives
\[
 v^{1-d}\le Kp^{1-d},\qquad v\le Kp.
\]
Dividing \Cref{a04:plank-lower} by the dummy-coordinate acceptance
bound gives coefficient mass at least
\(K^{-1}p^{d-1}v\ge p^d/K\).
The coefficient ball has radius \(Kp\) and meets only \(K\) ordinary
\(p\)-bins; one of them has the asserted mass.
All pruning and estimates used only the upper hypotheses and a dense
lower mass, proving the residual assertion.
\end{proof}

\subsection{Scalar-normal data and conditioning}

The next definition separates the fixed trajectory prior from the time
labels that we will repeatedly restrict.  Its three inputs control
scalar supports, comparisons of the row and residual slope, and spatial
mass conditional on a coefficient bin.

\begin{definition}[Scalar-normal data]\label{a04:scalar-data}
Fix \(j\in\{1,2\}\), \(0\le d<1\), and a floor \(\zeta\).
Let \(\nu\) be a trajectory probability, let \(x_l\) be quadratic and
\(y_l:[0,1]\to\R^j\) affine, and suppose their coefficients are bounded
by \(K\).  Let \(E_l\) have dense mass under \(\nu\otimes dt\), and let
\(F(t)\in(\R^j)^*\) be a measurable, time-only row of norm at most \(K\).
Write
\[
 S_l(t)=x_l(t)-F(t)y_l(t).
\]
For a bounded set \(A\), let \(N_q(A)\) count the cells of a fixed
half-open mesh of side \(q\) that meet \(A\), with the usual dyadic
rounding. The following bounds are required at the available resolutions:
\begin{enumerate}
\item For every active time, every interval \(J\) of length \(r\), and
\(\zeta\le q\le r\le1\),
\begin{equation}\label{a04:scalar-support}
 N_q\bigl(J\cap\{S_l(t):l\text{ with }t\in E_l\}\bigr)\le K(r/q)^d.
\end{equation}
For every spatial cube \(Q\subset\R^{1+j}\) of side \(q\),
\begin{equation}\label{a04:absolute-cap}
 \nu\{l:t\in E_l,\ (x_l(t),y_l(t))\in Q\}\le Kq^{j+d}.
\end{equation}
\item At active times,
\begin{equation}\label{a04:row-comparison}
 |F(t)-F(u)|\le K(|t-u|+\zeta).
\end{equation}
Writing \(R_l(t)=x_l'(t)-F(t)y_l'\), labeled spacetime points satisfy
\begin{equation}\label{a04:residual-comparison}
 |R_l(t)-R_{\tilde l}(u)|
 \le K\bigl(\dist((t,x_l(t),y_l(t)),
                  (u,x_{\tilde l}(u),y_{\tilde l}(u)))+\zeta\bigr).
\end{equation}
\item Let \(l_r\) bin all the displayed polynomial coefficients at side
\(r\), and let \(B\) be a participating bin of positive prior mass.
For every spatial \(q\)-cube, \(\zeta\le q\le r\),
\begin{equation}\label{a04:bin-cap}
 \frac{\nu\{l\in B:t\in E_l,\ (x_l(t),y_l(t))\in Q\}}{\nu(B)}
 \le K(q/r)^{j+d}.
\end{equation}
\end{enumerate}
These are absolute labeled masses relative to the stated prior.
A common null set of times may be removed.  A finite world parameter
is permitted, with the hypotheses holding in each retained conditional
world; every conditioning variable below then includes that world.
\end{definition}

We use fixed dyadic grids, allowing bounded enlargements.
All exponents of resolutions used in one application range over a fixed
finite interval.  Slowly densifying exponent grids give intermediate
covering bounds by subdivision, since consecutive scale ratios are
subpower.  The conditional cap \Cref{a04:bin-cap} is needed on compatible
parameter grids at the entropy scales.  A restriction of labels preserves
every upper hypothesis in \Cref{a04:scalar-data} with the same prior.

For these data we will approximate \(F\) by polynomials on labeled
portions of individual trajectories within short time blocks.  Each
portion will have inverse-subpower relative length and contain the
incidence being fitted; the polynomial may depend on that incidence.
The next entropy estimates express the obstruction behind this fit:
a scalar with support exponent \(d<1\) cannot determine a tangential
projection whose entropy has exponent one.  Later discrepancies will
be tested on events involving several times, so we prove these estimates
anew on every dense tested event.

For clarity, write \(\mathsf H\) and \(\mathsf I\) for Shannon entropy
and conditional mutual information, and set \(o_*=o(\log N)\).
Exact data are conditioned upon before only the displayed outputs are
discretized.

\begin{lemma}[Entropy on a dense tested event]\label{a04:dense-entropy}
Presample a finite time tuple independently of the trajectory in its
world, with absolutely continuous marginal laws for its tested times.
Let \(T\) contain the world and the entire tuple.
Let \(\mathcal A\) be any dense event on which all tested labels hold,
and normalize its restriction to a probability \(Q\).
At every tested time \(t\) and available dyadic scales,
\begin{align}
 \mathsf H_Q(S(t)_r\mid T)
   &=d\log(1/r)+o_*,\label{a04:scalar-entropy}\\
 \mathsf H_Q(y(t)_r\mid T,S(t)_r)
   &=j\log(1/r)+o_*,\label{a04:tangential-entropy}\\
 \mathsf I_Q(l_r;S(t)_q\mid T,S(t)_r)
   &=o_*,\qquad q<r.\label{a04:entropy-independence}
\end{align}
Moreover, for every \(T\)-measurable unit row \(e\) and \(r\ge q\),
\begin{equation}\label{a04:projection-entropy}
 \mathsf H_Q((e\,y(t))_r\mid T,S(t)_q)
       \ge\log(1/r)-o_*.
\end{equation}
The same assertions hold anew on every further dense event; the event
may depend on all the times and on the trajectory.
\end{lemma}

\begin{proof}
Let \(P\) be the law before restriction, \(p_{\mathcal A}=P(\mathcal A)\),
and \(s(T)=P(\mathcal A\mid T)\).  The following elementary denominator
bound will also be used with additional coefficient-bin conditioning:
\begin{equation}\label{a04:denominator-entropy}
 \E_Q\frac1{s(T)}\le\frac1{p_{\mathcal A}},
 \qquad
 \E_Q\log\frac1{s(T)}\le\log\frac1{p_{\mathcal A}}=o_*.
\end{equation}
The first inequality follows by integrating \(s(T)/p_{\mathcal A}\)
against the law of \(T\) on \(\{s>0\}\); the second is Jensen's inequality.
Thus no uniform lower bound for individual success probabilities is
required.

The shear \((x,y)\mapsto(S,y)\) is \(T\)-known and has subpower distortion.
Put \(Z=(S_l(t)_r,y_l(t)_r)\).  Since \(\mathcal A\) implies the
tested label at \(t\), every joint bin \(z\) satisfies
\[
 P(\mathcal A,Z=z\mid T)
 \le P(t\in E_l,Z=z\mid T)\le Kr^{j+d}.
\]
For the last inequality, the trajectory law at fixed \(T\) is the
original within-world prior, and the sheared bin is covered by \(K\)
spatial \(r\)-cubes to which \Cref{a04:absolute-cap} applies.
Thus, whenever \(s(T)>0\),
\[
 Q(Z=z\mid T)=\frac{P(\mathcal A,Z=z\mid T)}{s(T)}
       \le\frac{Kr^{j+d}}{s(T)}.
\]
The entropy is therefore at least
\[
 (j+d)\log(1/r)-O(\log K)-\E_Q\log(1/s(T)).
\]
The support bounds are \(Kr^{-d}\) for \(S_r\) and \(Kr^{-j}\) for \(y_r\).
Subtracting either support entropy proves
\Cref{a04:scalar-entropy,a04:tangential-entropy}, including their
matching upper bounds.

For \Cref{a04:entropy-independence}, put \(B=l_r\) and
\(s(T,B)=P(\mathcal A\mid T,B)\).
Before restriction the conditional trajectory law is exactly the
within-bin prior in \Cref{a04:bin-cap}, because the tuple was
trajectory-independent.  In \(B\), there are at most
\(K(r/q)^j\) relevant \(y_q\)-bins for a fixed \(S_q\)-bin.
Hence
\[
 P(\mathcal A,\ S_q=z\mid T,B)\le K(q/r)^d.
\]
The calculation in \Cref{a04:denominator-entropy} gives
\[
 \E_Q\log(1/s(T,B))\le\log(1/p_{\mathcal A}),
\]
with no dependence on the number or individual prior masses of \(B\).
Consequently
\[
 \mathsf H_Q(S_q\mid T,B)\ge d\log(r/q)-o_*.
\]
Since the coefficients vary by \(O(r)\) within \(B\), the time-known
shear permits only \(K\) values of \(S_r\) there.  Chain rule gives
\[
 \mathsf H_Q(S_q\mid T,B,S_r)
 \ge \mathsf H_Q(S_q\mid T,B)-\mathsf H_Q(S_r\mid T,B)
 \ge d\log(r/q)-o_*.
\]
On the other hand, \Cref{a04:scalar-support} bounds
\(\mathsf H_Q(S_q\mid T,S_r)\) above by \(d\log(r/q)+o_*\).
Their difference proves \Cref{a04:entropy-independence}.

For \Cref{a04:projection-entropy}, fix \(T\) and a projection bin of
width \(r\).  The corresponding slab in the bounded \(y\)-range meets
at most \(Krq^{-j}\) cubes of side \(q\).
Subtract its logarithm from \Cref{a04:tangential-entropy} at scale \(q\).
All estimates used the new event's actual probability
\(p_{\mathcal A}\); they can therefore be repeated on any dense event.
Absolutely continuous time marginals suffice to avoid the common null
set, and a density ceiling for the tuple law is unnecessary.
\end{proof}

We shall repeatedly use the following exact entropy inequality, valid
for any discrete \(Z,A,B,C\), with arbitrary further conditioning \(T\):
\begin{equation}\label{a04:common-information}
 \mathsf H(Z\mid T,C)
 \le \mathsf H(Z\mid T,B)+\mathsf H(Z\mid T,A)
          +\mathsf I(A;B\mid T,C).
\end{equation}
Indeed, expand the left side as
\(\mathsf H(Z\mid T,B,C)+\mathsf I(Z;B\mid T,C)\), enlarge \(Z\) to
\((A,Z)\) in the mutual information, and use chain rule.
In our applications, \(Z\) has negligible entropy both given
\((T,S_q)\) and given \((T,l_r)\).  Taking \(A=S_q\), \(B=l_r\),
and \(C=S_r\), the mutual-information estimate
\Cref{a04:entropy-independence} then gives negligible entropy
conditional on \((T,S_r)\) alone.  Boundary ambiguities caused by
bounded enlargements cost only \(O(\log K)\).

\subsection{The first jet}

\begin{proposition}[Scalar-normal first jet]\label{prop:scalar-normal}
Assume \Cref{a04:scalar-data}.  Fix a finite nested hierarchy of dyadic
time lengths of fixed positive power, with finest length \(H_*\).
If the available floor is sufficiently fine---\(\zeta\le H_*^3/K\)
suffices after enlarging \(K\)---then a dense restriction of the
single-time labels admits a time-only row \(A(t)\), bounded by \(K\),
such that
\begin{equation}\label{a04:first-taylor}
 |F(v)-F(u)-(v-u)A(u)|\le K H^{11/10}
\end{equation}
whenever the retained active times \(u,v\) lie in the same designated
interval of length \(H\).  Put
\[
 J_1(l,t)=x_l'(t)-F(t)y_l'-A(t)y_l(t).
\]
For every finite presampled tuple and every dense tested event as in
\Cref{a04:dense-entropy}, on which all tested times carry the retained
labels,
\begin{equation}\label{a04:first-jet-entropy}
 \mathsf H(J_1(l,t)_r\mid T,S_l(t)_r)=o_*
 \qquad\text{if }r\ge H_*^{(1-d)/64}
\end{equation}
at the available resolutions.  The row \(A\) and the single-time
restriction are fixed before the tuple and the later event are chosen.
\end{proposition}

The first step is geometric: we fit each component of \(F\) by an
affine function on selected times in each hierarchy block.  After local
normalization, the residual comparison expresses the residual slope as an
approximately Lipschitz function of one tangential coordinate.
Collisions of scalar values relate this graph to the graph of the
corresponding component of \(F\) through dot products.  The scalar
support has exponent \(d<1\), which will exclude the expansion
alternative in the following theorem.  Its line-concentration
alternative then supplies the affine fit.

Comparing these fits across lengths will produce the row \(A\) and
its Taylor estimate.  We will then concentrate the residual slopes
\(J_1\).  This last step gives fixed scalar and coefficient-bin covers,
so that the entropy estimate can be proved on every later dense tested
event.


\begin{theorem}[Planar dot-product alternative]\label{a04:planar-input}
For every \(\varepsilon>0\) there are \(\eta,\delta_0>0\) such that the
following holds for \(0<p\le\delta_0\).
Let \(\mathcal A,\mathcal B\subset[0,1]^2\), and suppose that
\begin{equation}\label{a04:planar-regularity}
 N_s(E\cap B(z,w))\le p^{-\eta}w/s
 \quad\text{for }E\in\{\mathcal A,\mathcal B\},\quad p\le s\le w.
\end{equation}
At least one of the following alternatives holds:
\begin{enumerate}
\item There are affine lines \(\ell,\ell^\perp\) with perpendicular
directions such that both
\[
 N_p(\mathcal A\cap\{z:\dist(z,\ell)\le p\}),\qquad
 N_p(\mathcal B\cap\{z:\dist(z,\ell^\perp)\le p\})
\]
are at least \(p^{\varepsilon-1}\).
\item Every restricted triple set
\(\mathcal H\subset\mathcal A\times\mathcal B\times\mathcal B\)
with \(N_p(\mathcal H)\ge p^{\eta-3}\) admits a scale \(s\ge p\) and
an interval \(J\), \(|J|\ge p^{-\eta}s\), for which
\[
 N_s\bigl(J\cap
       \{a\cdot(b-b'):(a,b,b')\in\mathcal H\}\bigr)
       \ge(|J|/s)^{1-\varepsilon}.
\]
\end{enumerate}
\end{theorem}

We use the restatement of Theorem~5.2 at the start of Section~8 of
\citet{WangZahlSticky}, together with Definition~4.2, with theorem
numbering fixed to arXiv:2210.09581v2 (2025).
In particular, \Cref{a04:planar-regularity} is only an upper covering
condition, and the second alternative concerns a restricted triple set.
We use the first alternative only for \(\mathcal A\).

\begin{proof}[Proof of \Cref{prop:scalar-normal}]
Before any label selection, insert finitely many intermediate dyadic
lengths so that adjacent lengths \(H'\le H\) satisfy
\(H'\ge H^{33/32}/2\).  They will permit comparison of the affine
slopes after the fits have been constructed.

We construct those fits one row component and one hierarchy length at
a time.  The row \(F\) is common to all trajectories.  A temporary
group of trajectories can therefore locate fitting times, after which
we may retain all current labels at those times.  The next preparation
ensures that this recovery retains dense incidence mass.

\paragraph{Preparing the current labels.}
We perform the following preparation before fitting each row component
at each hierarchy length \(H\).  Denote the current labels by
\(E_l^{\rm cur}\), and write
\[
 m(t)=\nu\{l:t\in E_l^{\rm cur}\},\qquad
 \delta=\int_0^1m(t)\,dt\ge K^{-1}.
\]
Choose a reciprocal-subpower threshold \(\tau=o(\delta)\).
First remove times with \(m(t)<\tau\).  Next remove each dyadic
\(H\)-block whose remaining incidence mass is less than \(\tau H\).
The two removals cost at most \(\tau\) each, so the resulting labels
\(E_l^{\rm pre}\) have mass at least \(\delta/2\).
Every surviving block has mass at least \(\tau H\ge H/K\), and at
every surviving active time the instantaneous mass is still
\(m(t)\ge\tau\): both removals kept or discarded all current labels
at a time.  Since \(\nu\) is a probability, each block has mass at
most \(H\).  Thus at least \(\delta/(2H)\) blocks survive.
We now run the local construction separately in each of these blocks.

\paragraph{Local normalization.}
Fix a surviving block \(I=[t_I,t_I+H]\), put \(u=(t-t_I)/H\), and define
\[
 U_l=\{u\in[0,1]:t_I+Hu\in E_l^{\rm pre}\cap I\}.
\]
Averaging and Cauchy--Schwarz provide \(u_*\) such that
\[
 \int \1_{\{u_*\in U_l\}}|U_l|\,d\nu(l)\ge K^{-1}.
\]
During the local construction retain trajectories labeled at
\(t_*=t_I+Hu_*\), and put
\[
 F_*=F(t_*),\qquad B_l=x_l(t_I)-F_*y_l(t_I).
\]
The difference \(B_l-S_l(t_*)\) is \(O(KH)\).
Thus the occupied \(H\)-bins \(C\) of \(B_l\) number at most \(KH^{-d}\),
and each has prior mass at most \(KH^d\).
For the latter statement, sum \Cref{a04:absolute-cap} at time \(t_*\)
over the \(KH^{-j}\) tangential \(H\)-bins compatible with
\(|B_l-S_l(t_*)|\le KH\).
In each \(C\), divide the restricted prior by \(H^d\) and call the
result \(\nu_C\).  Deleting groups with
\(\int |U_l|\,d\nu_C<K_1^{-1}\) loses at most \(K/K_1\) of the original
normalized incidence mass.  Choose a sufficiently large subpower
\(K_1\), so that a dense family of good groups remains.

For a lower endpoint \(b_C\), define
\begin{align*}
 Y_l&=y_l(t_I),&
 z_l&=(B_l-b_C)/H,&
 q_l&=x_l'(t_I)-F_*y_l',\\
 \phi(u)&=(F(t_I+Hu)-F_*)/H,&
 L_l(u)&=z_l+u q_l-\phi(u)Y_l.
\end{align*}
The normalized row \(\phi\) measures the variation that we must fit
on this block.  The scalar model \(L_l\) separates that row from the
frozen position \(Y_l\) and residual slope \(q_l\).
These quantities are bounded by \(K\) on the labels, and \(\phi\) obeys
a Lipschitz comparison with slack \(\zeta/H\).
Writing \(x_{2,l}\) for the quadratic coefficient, direct expansion gives
the exact identity
\begin{equation}\label{a04:quadratic-local-error}
 \frac{S_l(t_I+Hu)-b_C}{H}
 =L_l(u)+H\bigl(u^2x_{2,l}-u\phi(u)y_l'\bigr).
\end{equation}
The final term is \(O(KH)\).
Consequently at \(p=H^{1/4}\), and also at \(p_0=H^{1/8}\), the
\(L_l(u)\)-supports retain the upper \(d\)-count.
Moreover,
\begin{equation}\label{a04:local-joint-cap}
 \nu_C\{l:u\in U_l,\ (Y_l,L_l(u))\text{ in given }p\text{-bins}\}
       \le Kp^{j+d}.
\end{equation}
Indeed, in the original variables the corresponding region has
tangential widths \(Kp\) and normal width \(KHp\).  It is covered by
\(KH^{-j}\) cubes of side \(Hp\), and division by \(H^d\) converts their
total cap to \(Kp^{j+d}\).
The same proof gives \Cref{a04:local-joint-cap} with \(p_0\).
Finally, the residual comparison at \(t_*\) gives
\begin{equation}\label{a04:frozen-q-comparison}
 |q_l-q_{\tilde l}|
       \le K(|Y_l-Y_{\tilde l}|+H).
\end{equation}
Here \(y_l(t_*)=Y_l+O(KH)\), the relevant scalar values differ by
\(O(KH)\), and replacing \(x_l'(t_*)\) by \(x_l'(t_I)\) costs \(O(KH)\).
Thus the quadratic part of \(x_l\) contributes only the recorded
\(O(KH)\) errors.  At the coarser scale \(p=H^{1/4}\), freezing the
other coordinates of \(Y_l\) makes \(q_l\) an approximately Lipschitz
function of the remaining coordinate.  We now use the planar theorem
to fit the corresponding component of \(\phi\).

\paragraph{Fitting one row component.}
Slice the other \(j-1\) coordinates of \(Y_l\) into \(p\)-bins; no slice
is needed when \(j=1\).
A group consisting of \(C\) and one such slice has prior mass at most
\(KH^dp^{j-1}\), by the spatial cap at \(t_*\).
After division by \(H^dp^{j-1}\) and deletion of light groups, choose a
sliced measure \(\nu_s\) with
\[
 \nu_s(\text{all trajectories})\le K,\qquad
 \int|U_l|\,d\nu_s(l)\ge K^{-1}.
\]
Let \(v_l\) be the remaining coordinate of \(Y_l\).
Choose representatives of \(q_l\) on occupied \(v\)-bins and interpolate
between their centers.  \Cref{a04:frozen-q-comparison} gives a function
\(G\) satisfying
\[
 |G(v)-G(v')|\le K(|v-v'|+p).
\]
Translate \(L_l(u)\) by the other components of \(\phi(u)\) times the
slice centers.  The resulting labeled values satisfy
\begin{equation}\label{a04:sliced-model}
 L_l^a(u)=z_l+uG(v_l)-\phi_1(u)v_l+O(Kp).
\end{equation}
They have the local upper \(d\)-count, and
\begin{equation}\label{a04:sliced-cap}
 \nu_s\{l:u\in U_l,\ (v_l,L_l^a(u))
             \text{ in given }p\text{-bins}\}\le Kp^{1+d}.
\end{equation}

Choose \(u_0\) with
\(\int\1_{\{u_0\in U_l\}}|U_l|\,d\nu_s\ge K^{-1}\).
For pairs of trajectories labeled at both \(u_0\) and \(u\), count
collisions of their \(L^a(u)\)-bins using the unnormalized measure
\(d\nu_s(l)\,d\nu_s(l')\,du\).
Cauchy--Schwarz and the \(Kp^{-d}\) output-bin count give collision
mass at least \(p^d/K\).
Fixing the first trajectory's starting scalar bin at \(u_0\) leaves
at least \(p^{2d}/K\) of this mass.
On the retained pairs, \Cref{a04:sliced-model} implies
\begin{align}
 &(u-u_0)(G(v_l)-G(v_{l'}))
       -(\phi_1(u)-\phi_1(u_0))(v_l-v_{l'})\notag\\
 &\hspace{35mm}=L_{l'}^a(u_0)-L_l^a(u_0)+O(Kp).
 \label{a04:collision-dot}
\end{align}
At fixed \(u\) and fixed \(p\)-bins of \(v_l,v_{l'}\), the first
trajectory has weight at most \(Kp^{1+d}\) by its fixed starting
scalar bin.  The equation confines the second starting scalar value
to \(K\) bins, so its weight has the same bound.
Each \((u,v_l,v_{l'})\)-cell therefore carries at most \(Kp^{3+2d}\).
The total pair weight at a fixed \(u\) is at most \(Kp^{2d}\):
sum \Cref{a04:sliced-cap} over the tangential bins for the first
starting value and the second matching value.

For each time \(p\)-bin \(J\), let \(T_J\subset J\) be its active
pair times for the retained collisions, before any quadrant restriction.
Discard bins with \(|T_J|<p/K_1\), and keep these pre-quadrant sets
\(T_J\) attached to all surviving bins.
There are \(O(p^{-1})\) time bins, so the discarded pair mass is at most
\(Kp^{2d}/K_1\).  Take \(K_1\) large enough to preserve the lower bound.
The represented triples
\[
 a=(u-u_0,-\phi_1(u)+\phi_1(u_0)),\quad
 b=(G(v_l),v_l),\quad b'=(G(v_{l'}),v_{l'})
\]
then have \(p\)-covering number at least \(p^{-3}/K\).

Retain a quadrant for \(a\) containing a fixed fraction of the pair
mass.  Let \(\mathcal H\) be the image of these retained collision
triples, let \(\mathcal A\) be its first-coordinate projection, and
let \(\mathcal B\) be the union of its two trajectory-coordinate
projections.  Thus
\(\mathcal H\subset\mathcal A\times\mathcal B\times\mathcal B\), with
the same lower covering count up to a fixed factor.
Reflect both planar sets by the same coordinate reflection,
dilate each by a positive scalar between \(K^{-1}\) and \(1\), and
translate only the second set to place the sets in \([0,1]^2\).
Translation of the second set cancels in \(b-b'\); no translation of
the first set is made.
Their graph comparisons give the upper covering condition
\Cref{a04:planar-regularity} with a subpower constant: covering a
horizontal interval for the first graph, or a vertical interval for
the second, by \(s\)-intervals gives at most \(Kw/s\) graph cells.
Normalization changes this constant and the triple lower count only
by subpowers.

Fix \(0<\varepsilon<1-d\).  For sufficiently large \(N\), the
\(p^{\pm\eta}\) allowances in \Cref{a04:planar-input} absorb these
subpower losses.
Its expansion alternative is impossible: by
\Cref{a04:collision-dot}, all dot products are within \(Kp\) of a
translate of the starting scalar support, so they use at most
\(K(|J|/s)^d\) \(s\)-bins in \(J\).
This would imply
\[
 (|J|/s)^{1-\varepsilon-d}\le K,
 \qquad |J|/s\ge p^{-\eta},
\]
a contradiction.
The additive error and inverse dilations cost only subpowers at
\(s\ge p\); intervals larger than the unit cutoff can be split into
subpower many intervals in the bounded range.

We therefore obtain at least \(p^{\varepsilon-1}/K\) time bins near one
line in the graph of \(\phi_1\).
The line is a graph over time with slope at most \(K\):
two represented points have horizontal separation at least
\(p^\varepsilon/K\), larger than their \(Kp\) proximity errors, and
the graph comparison bounds their slope.
Let \(\mathcal J_I\) be the resulting family of represented time
bins.  In each \(J\in\mathcal J_I\), the graph comparison extends the
affine fit from a represented near-line point to every time in the
stored set \(T_J\), with error \(Kp\).  Indeed, their time separation
is at most \(p\), and both the graph comparison and the fitted line
have slope bounded by \(K\).  Use precisely the time set
\[
 T_I=t_I+H\bigcup_{J\in\mathcal J_I}T_J.
\]
Every \(T_J\) has length at least \(p/K_1\); the disjoint bins and
their number therefore give
\[
 |T_I|\ge Hp^\varepsilon/K.
\]
In particular, the fitting set consists of prepared active times,
not the whole selected bins.
Let \(\varepsilon\) tend to zero sufficiently slowly, after the
corresponding \(\eta,\delta_0\) and large-\(N\) thresholds have been
fixed.  Then \(|T_I|\ge H/K\), while the normalized fit error remains
\(Kp\).

In the original variables this is an affine component fit to \(F\)
with error \(KH^{5/4}\) on \(T_I\).
At these times retain \emph{all} prepared labels \(E_l^{\rm pre}\),
including those outside the temporary sliced group.  Their total
incidence mass satisfies
\begin{equation}\label{a04:fit-mass-recovery}
 \sum_{I\ \mathrm{surviving}}\int_{T_I}m(t)\,dt
 \ge\tau\sum_I|T_I|
 \ge\tau\frac{\delta}{2H}\frac{H}{K}
 =\frac{\delta\tau}{2K}.
\end{equation}
This is inverse-subpower.  It uses both the instantaneous floor and
the number of surviving blocks, each established before the local
search.  For \(j=2\), repeat with the second component on these new
labels, recomputing \(m,\delta,\tau\) first.  Do the same at every
length in the fixed finite hierarchy.  Subsequent label restrictions
preserve the earlier fits and all upper inputs.  We obtain affine rows
\(P_I(t)\), with slopes
\(A_I\) of size at most \(K\), satisfying
\begin{equation}\label{a04:vector-affine-fit}
 |F(t)-P_I(t)|\le KH^{5/4}
\end{equation}
on every retained active time in the corresponding interval.
Let \(I_*(t)\) be the finest dyadic block containing \(t\), using a
fixed half-open endpoint convention.  Define
\[
 A(t)=A_{I_*(t)}
\]
on the finest blocks with retained labels, and set \(A(t)=0\)
elsewhere.  This is a fixed time-only row bounded by \(K\).
Subsequent pruning restricts labels without changing this row.

\paragraph{Comparison across lengths.}
We now compare the affine fits through the intermediate lengths fixed
before any selections.  Let \(\mu_{\rm fit}\) be the current labeled
measure.  At every hierarchy length \(H\), delete the labels in
blocks \(I\) with \(\mu_{\rm fit}(I)<H/K_1\), testing all blocks
against this same measure.  Each length costs at most \(1/K_1\), so
choose \(K_1\) large enough to make the total loss negligible.
Every child block that still contains a retained incidence had
\(\mu_{\rm fit}\)-mass at least \(H'/K_1\).  Since \(\nu\) is a
probability, its fitting times have length at least \(H'/K_1\).
Both its own affine fit and its parent's hold on those times,
whether or not all of them survive the deletion.  Choose two separated
by \(H'/K\).  Comparing the fits there gives
\[
 |A_{I'}-A_I|
 \le K\frac{H^{5/4}}{H'}\le KH^{7/32}.
\]
Summing through the fixed finite chain compares the finest slope
\(A(u)\) with \(A_I\) by \(KH^{7/32}\).
Together with \Cref{a04:vector-affine-fit}, this gives
\[
 |F(v)-F(u)-(v-u)A(u)|
 \le K\bigl(H^{5/4}+H^{1+7/32}\bigr)\le KH^{11/10}.
\]
This proves \Cref{a04:first-taylor}, which persists under the label
restrictions below.

\paragraph{Concentrating the residual slopes.}
The affine fits now hold on a dense set of incidences.  We next restrict
the trajectory labels so that their residual slopes admit short lists
once the time and scalar value are known.
At the finest length \(H=H_*\), repeat the unsliced \(C\)-group
normalization on the retained labels.
With \(p_0=H^{1/8}\), the affine approximation to \(\phi\) gives an
affine approximation to \(L_l\) with slope
\[
 k_l=q_l-A_IY_l
\]
and error at most \(Kp_0\).
The slopes vary by at most \(Kp_0\) in each \(Y\)-cell of side \(p_0\),
by \Cref{a04:frozen-q-comparison}.
Put \(\gamma=(1-d)/4\).
We claim that some interval of width \(p_0^\gamma\) contains slopes of
labeled incidence mass at least \(K^{-1}\) in every good group,
with uniform subpower constants.

If this failed, there would be a sequence of groups and a fixed
\(c>0\) on which every such interval has incidence mass at most
\(N^{-c}\).  Let \(\tau=N^{-c/4}\).
Delete trajectories with labeled length below \(\tau\), and then
\(Y\)-cells whose remaining trajectory mass is below \(p_0^j\tau\).
The two losses are at most \(K\tau\).
The remaining output collisions, integrated in time, still have mass
at least \(p_0^d/K\), by Cauchy--Schwarz.

For a fixed first trajectory, second trajectories with slope within
\(p_0^\gamma\) occupy at most
\[
 K N^{-c}\tau^{-2}p_0^{-j}
\]
retained \(Y\)-cells.
Indeed, their slopes, including the \(Kp_0\) variation within each cell,
lie in a bounded number of intervals of the assumed small incidence
mass.  Their labeled lengths are at least \(\tau\), so their prior
mass is at most \(KN^{-c}/\tau\); each retained cell has prior mass
at least \(p_0^j\tau\).
For each such cell the matching output weight at a time is at most
\(Kp_0^{j+d}\), by \Cref{a04:local-joint-cap}.
Thus the close-slope collision contribution is at most
\[
 Kp_0^d N^{-c}\tau^{-2}=Kp_0^dN^{-c/2}.
\]
For the other pairs, the two affine approximations can agree to
\(Kp_0\) only on times of length at most \(Kp_0^{1-\gamma}\).
Their total contribution is at most that quantity times a subpower
bound for the product prior mass.
Both contributions are \(o(p_0^d/K)\), because \(1-\gamma>d\).
This contradicts the collision lower bound and proves the claim.
If the claim's subpower constants were not uniform across groups,
an offending sequence of groups would give exactly the contradiction
just proved.

Keep one such slope interval in each good group.
Since
\begin{equation}\label{a04:quadratic-jet-error}
 J_1(l,t)-k_l
   =H\bigl(2u x_{2,l}-\phi(u)y_l'-uA_Iy_l'\bigr)=O(KH),
\end{equation}
this gives a deterministic interval cover for \(J_1\).
For fixed \(t,S_l(t)_H\), only \(K\) base bins \(C\) are possible:
\(|B_l-S_l(t)|\le KH\).
Hence \(J_1\) lies in \(K\) intervals of width
\[
 Kp_0^\gamma=KH^{(1-d)/32}.
\]
At any reading scale \(r\ge H^{(1-d)/64}\), these are \(K\) \(r\)-bins.
The cover survives appending all tuple times to \(T\) and imposing
any later event.  The formula for \(J_1\), whose rows are now fixed
and bounded, also gives \(K\) \(r\)-bins when \(T,l_r\) is fixed.
Apply \Cref{a04:common-information} and
\Cref{a04:entropy-independence} anew on each dense tested event,
with fine scalar \(S_H\) and coarse scalar \(S_r\).
This proves \Cref{a04:first-jet-entropy}.
The argument rederives the entropy estimate on the specified event,
using covers that persist under restriction.  Together with the Taylor
estimate already proved, this completes the proposition.
\end{proof}


\subsection{Higher jets and polynomial fits}

Fix \(0\le d<1\), an integer \(M\ge2\), and a dyadic block length
\(h=N^{-c+o(1)}\), where \(c>0\) is fixed.  We seek a degree-\(M\)
polynomial fitting the row on a labeled portion of a trajectory of
length at least \(h/K\).  The first-jet
information will be differentiated along a finite tree of nearby times.
At each pair scale \(H\), the child jets must already be readable at
width \(H^{1.01}\), so that division by the pair gap still leaves a
positive power of accuracy.  The finest scale supplies the first jets;
the root accuracy must be finer than the error used in the final
polynomial test.  We fix all these scales before selecting labels.

Define
\[
 c_0=.1,\qquad
 r_i=\min(.009,c_{i-1}/2),\qquad c_i=r_i/2
 \quad(1\le i\le M-1),
\]
and put
\[
 r_{\min}=\min_{1\le i\le M-1}r_i
          =.009\,4^{-(M-2)},\qquad
 \alpha_0=(1-d)/64.
\]
Both numbers are positive and depend only on \(M,d\).
Starting with \(h\), choose nested dyadic lengths
\[
 h\gg H_1\gg H_2\gg\cdots\gg H_{M-1}\gg H_*,
\]
by taking at each step the largest dyadic length satisfying the
corresponding inequality
\begin{equation}\label{a04:jet-scale-schedule}
 \begin{split}
 H_1^{r_{\min}}&\le h^{M+11},\\
 H_i^{r_{\min}}&\le H_{i-1}^{1.02}
                  \quad(2\le i\le M-1),\\
 H_*^{\alpha_0}&\le H_{M-1}^{1.02}.
 \end{split}
\end{equation}
Choose the available floor to satisfy
\begin{equation}\label{a04:jet-floor}
 \zeta\le K^{-1}\min\{H_*^3,h^{M+1/2}\}.
\end{equation}
These choices use finite fixed power ratios.  We will verify the
required entropy readings in the proof.  The strict exponent margins
absorb dyadic rounding and subpower factors.  When specifying the floor
by a fixed power, one may place it a further fixed power of \(h\) below
the displayed minimum to allow every later subpower enlargement of
\(K\).

\begin{proposition}[Higher-jet polynomial fitting]\label{prop:higher-jets}
For the parameters \(M,d,h\) and hierarchy just defined, assume
\Cref{a04:scalar-data}, including the parameter-bin cap, and the floor
condition \Cref{a04:jet-floor}.  Then a dense set of labeled incidences
\((l,t)\) has the following property.
There is a row-valued polynomial \(P\) of degree at most \(M\), with
coefficients bounded by \(K\), such that in the \(h\)-block containing
\(t\),
\[
 |F(v)-P(v)|\le Kh^{M+1/4}
\]
on labeled times \(v\in E_l\) of length at least \(h/K\), including \(v=t\).
The polynomial may depend on the incidence.
The conclusion applies anew to any dense restriction of the input
labels with the same upper hypotheses.
\end{proposition}

\begin{proof}
Use the hierarchy and floor fixed in
\Cref{a04:jet-scale-schedule,a04:jet-floor}.
Apply \Cref{prop:scalar-normal}, inserting its finite intermediate
Taylor scales before making any label selections.
We retain the resulting single-time labels and the row \(A(t)\).

\paragraph{An independent time tree.}
Choose \(u\) uniformly in \([0,1]\).
Add a second time uniformly in its \(H_1\)-interval.
At the next level, split each of these times into itself and a new
independent uniform draw in its \(H_2\)-interval.
Continue for \(M-1\) pair levels.
Separately draw \(t'\) uniformly in the \(h\)-block of \(u\), independently
of all the other draws conditional on that block.
This entire experiment is independent of the trajectory before any
labels are imposed.  Let \(T\) contain the world, the whole tree, and
\(t'\).

All tested times are labeled with dense probability.
Here is a direct verification.  For a fixed trajectory and a fixed
containing interval, write \(a\) for the fraction of labeled times.
At a split, conditional on the smaller partition interval, the two
child times are independent uniform points of that interval.
The mean of the product of their subtree success probabilities is
the square of the mean subtree probability.
Jensen's inequality, repeated up the finite tree, gives a lower bound
by \(a^{2^{M-1}}\) in each \(h\)-block.
The additional independent \(t'\) multiplies this by \(a\).
Averaging over blocks, trajectories and worlds, and using Jensen once
more, bounds total success below by the
\((2^{M-1}+1)\)-st power of the single-time label mass.
This is inverse-subpower.

We may also require that every pair gap at scale \(H_i\) be at least
\(H_i/K\).  Before label restriction, each excluded gap has probability
at most \(2/K\).  Since the number of pairs is fixed, take this \(K\)
large enough compared with the reciprocal all-label probability.
The loss is then negligible relative to that probability.
These restrictions are performed on the presampled tuple law.
We shall reapply \Cref{a04:dense-entropy} on each further dense event;
we do not assume independence after restriction.

\paragraph{Inductive information statement.}
Put \(F_0=F\) and \(F_1=A\).
Starting at the leaves, a subtree rooted at time \(t\) will carry rows
\(F_1(t),\ldots,F_m(t)\) and jets
\begin{equation}\label{a04:jet-definition}
 J_i(t)=x_l^{(i)}(t)-F_i(t)y_l(t)-iF_{i-1}(t)y_l',
 \qquad 1\le i\le m.
\end{equation}
Every row is a function of its subtree times and the world only.
In particular it is independent of the trajectory and does not use
the separate time \(t'\).  The rows are bounded by \(K\) on the retained
tuples.
At the readings required below, we maintain
\begin{equation}\label{a04:hereditary-jets}
 \mathsf H(J_i(t)_r\mid T,S_l(t)_r)=o_*
\end{equation}
on \emph{each} further dense event in the original trajectory--tuple
experiment.  At the leaves this is
\Cref{a04:first-jet-entropy} for the fixed row \(A\).
Its proof remains valid with the whole tree and \(t'\) appended to
\(T\): the deterministic fine-scalar and coefficient-bin covers
survive, and \Cref{a04:dense-entropy} is proved anew on the chosen event.

Consider a pair \(t,w\) at scale \(H\), whose two subtrees have jets
through order \(m\).  Set \(q=H^{1.01}\), with dyadic rounding.
Their internal scales will have been chosen fine enough for the
child assertions at \(q\).
Taylor expansion of \(x_l,y_l\), and \Cref{a04:first-taylor}, give
\begin{equation}\label{a04:scalar-first-expansion}
 S_l(w)=S_l(t)+(w-t)J_1(t)+O(KH^{1.1}).
\end{equation}
The error is smaller than \(q\).
Thus \(S_l(w)_q\), and then all the \(q\)-jets at both endpoints, have
negligible entropy given \(T,S_l(t)_q\).
This uses the child information statements on the current event and
chain rule.  Write
\[
 \nabla Z=\frac{Z(w)-Z(t)}{w-t},
\]
where endpoint rows use their own subtree data.
Since \(|w-t|\ge H/K\), the \(\nabla J_i\) are determined to accuracy
\(KH^{.01}\), with entropy cost \(o_*\), by \(T,S_l(t)_q\).

For \(i\ge1\), quadraticity of \(x_l\) and affinity of \(y_l\) give the
exact identity
\begin{equation}\label{a04:jet-difference}
 \nabla J_i
 =x_l^{(i+1)}(t)-(\nabla F_i)y_l(t)
       -\bigl(F_i(w)+i\nabla F_{i-1}\bigr)y_l'.
\end{equation}
For \(i=1\), the difference quotient of \(x_l'\) equals \(x_l''\);
for \(i\ge2\), both sides of the corresponding derivative identity
are zero.  The remaining terms in \Cref{a04:jet-difference} are simply
the exact product difference for an affine \(y_l\).
There is therefore no unrecorded higher-order scalar remainder.

Proceed through \(i=1,\ldots,m\), using the exponents fixed above and
maintaining
\begin{equation}\label{a04:row-compatibility}
 \nabla F_{i-1}=F_i(t)+O(KH^{c_{i-1}}),
 \qquad c_{i-1}>0.
\end{equation}
Initially \Cref{a04:first-taylor}, divided by the pair gap, supplies
the chosen exponent \(c_0=.1\).

First, \(\nabla F_i\) is bounded by a subpower after negligible loss.
If a fixed-power tail \(|\nabla F_i|\ge N^\alpha\), \(\alpha>0\), had
dense mass on a subsequence, divide \Cref{a04:jet-difference} by that
norm.
The other displayed coefficients are bounded by \(K\), using
\Cref{a04:row-compatibility}, and \(x_l^{(i+1)},y_l'\) have size \(K\).
The equation would therefore encode the projection of \(y_l(t)\)
in the \(T\)-known direction
\((\nabla F_i)/|\nabla F_i|\), with negligible entropy, at some fixed
positive power accuracy coarser than \(q\).
For example, if \(H=N^{-b+o(1)}\), take a reading
\(N^{-\beta}\) with \(0<\beta<\min(\alpha,b)\).
All nuisance terms after division are below this reading, as is the
divided \(q\)-bin error.
This contradicts \Cref{a04:projection-entropy} on that very tail event.
It follows that every fixed-power tail is negligible relative to
the current dense tuple mass, and a slow choice of cutoffs yields
\(|\nabla F_i|\le K\).

We now also have \(F_i(w)-F_i(t)=O(KH)\).
Recall that
\[
 r_i=\min(.009,c_{i-1}/2)>0.
\]
\Cref{a04:jet-difference,a04:row-compatibility} encode
\begin{equation}\label{a04:next-jet}
 U_{i+1}
 =x_l^{(i+1)}(t)-(\nabla F_i)y_l(t)
                      -(i+1)F_i(t)y_l'
\end{equation}
to width \(H^{r_i}\), at entropy cost \(o_*\), given \(T,S_l(t)_q\).
Indeed, the error \(KH^{.01}\), the \(O(KH)\) row difference, and the
\(O(KH^{c_{i-1}})\) compatibility error are all smaller than this width.

If \(i<m\), subtract the already known jet \(J_{i+1}(t)\).
This encodes
\((\nabla F_i-F_{i+1}(t))y_l(t)\) to the same width.
On an event where the row norm is at least \(H^{r_i/2}\), division would
encode its unit projection at width \(H^{r_i/3}\), again contradicting
\Cref{a04:projection-entropy} on any dense such event.
After a negligible deletion,
\[
 \nabla F_i-F_{i+1}(t)=O(KH^{r_i/2}).
\]
This proves the compatibility assertion with the chosen
\(c_i=r_i/2\), allowing us to proceed to the next value of \(i\).

For \(i=m\), define \(F_{m+1}(t)=\nabla F_m\) and keep the older rows
at \(t\).  This row is computed only from the two child rows and their
times; no trajectory or \(t'\)-dependent choice is made.
The new jet equals \(U_{m+1}\).
Its bin at \(r=H^{r_m}\), and at coarser required powers, has negligible
entropy given \(T,S_l(t)_q\).
Given \(T,l_r\), the bounded rows and coefficient-bin diameter allow
only \(K\) bins for that jet.
Applying \Cref{a04:common-information} with
\[
 Z=J_{m+1}(t)_r,\quad A=S_l(t)_q,\quad B=l_r,\quad C=S_l(t)_r
\]
and invoking \Cref{a04:entropy-independence} on the current event gives
\Cref{a04:hereditary-jets} for the new jet.
The older jets retain their hereditary assertions.

There are finitely many rows and discrepancy tests for fixed \(M\).
For each fixed positive threshold exponent, a nonnegligible relative
bad mass would itself yield a dense event on a subsequence, contradicting
the projection estimate just used.
Choose threshold exponents decreasing to zero so slowly that the
union of the finitely many bad sets has relative mass tending to zero.
The resulting row bounds are deterministic on the retained tuples.
Every later dense event inherits those deterministic bounds; its
entropy assertions are then derived anew from the child assertions,
\Cref{a04:common-information}, and
\Cref{a04:dense-entropy}.
We never infer hereditary entropy by conditioning an earlier average.
These cutoff tests concern rows already defined from the time tree.
They do not introduce a dependence on the trajectory or on \(t'\).

We can now verify the schedule.  A jet formed at a pair scale \(H_i\)
is available at width \(H_i^{r_m}\) for its relevant order \(m\), and
at all coarser required readings.  Since \(r_m\ge r_{\min}\), the
root widths satisfy
\[
 KH_1^{r_m}\le KH_1^{r_{\min}}\le Kh^{M+11}
          \ll p:=h^{M+10}.
\]
At a child scale \(H_i\), \(i\ge2\), the corresponding bound is
\[
 KH_i^{r_m}\le KH_{i-1}^{1.02}\ll H_{i-1}^{1.01},
\]
which supplies the reading \(q\) requested by its parent.  Older jets
inherited from finer subtrees obey the same requirements.  For the
leaf jets, \Cref{a04:first-jet-entropy} applies at every reading at
least \(H_*^{\alpha_0}\); by \Cref{a04:jet-scale-schedule}, this is
finer than \(H_{M-1}^{1.01}\), with the same fixed margin.
Finally, \(\zeta\le H_*^3/K\) covers the first-jet construction, its
fine scalar scale \(H_*\), and all the intermediate readings just
used.  The projection readings in the tail and compatibility tests
are coarser than their scalar conditioning scale \(q\).
Thus the entire induction uses available scales.

\paragraph{Testing the polynomial at a separate time.}
At the root \(u\), set
\[
 P_F(v)=\sum_{i=0}^M\frac{F_i(u)}{i!}(v-u)^i,\qquad
 E=F(t')-P_F(t'),\qquad J_0(u)=S_l(u).
\]
The coefficients of \(P_F\) are bounded by \(K\), including when
expanded in ordinary powers of \(v\), because \(M\) is fixed and
\(u\in[0,1]\).
Crucially, the coefficients are functions of the world and the tree
excluding \(t'\).
Exact evolution of the quadratic \(x_l\) and affine \(y_l\) gives
\begin{equation}\label{a04:final-jet-expansion}
 S_l(t')
 =\sum_{i=0}^M\frac{J_i(u)}{i!}(t'-u)^i
       -E\,y_l(t')+O(Kh^{M+1}).
\end{equation}
The last term is the single extra degree produced when the degree
\(M\) polynomial \(P_F\) multiplies the affine \(y_l\).

Suppose \(|E|\ge h^{M+1/4}\) on a dense event.
Work on that event and condition on \(T,S_l(u)_p\).
The root jet bins at \(p\) have total additional entropy \(o_*\);
after specifying them, \Cref{a04:final-jet-expansion} places \(S_l(t')\)
in an interval of length \(K|E|\).
Put \(r=h^{1/4}\).
The local support count \Cref{a04:scalar-support} allows at most
\(Kr^{-d}\) values at width \(|E|r\) in that interval.
Both the width and interval length are \(T\)-known.
The variable width satisfies
\[
 |E|r\ge h^{M+1/4}h^{1/4}=h^{M+1/2}\ge K\zeta,
\]
by \Cref{a04:jet-floor}; its containing interval is larger still.
Thus both are within the available scale range;
dyadic enlargement or the unit cutoff costs only \(K\).

After division by \(|E|\), the Taylor remainder is \(O(Kh^{3/4})\);
the root-jet bin error is smaller, and the scalar-bin error is \(O(r)\).
Also \(y_l(t')=y_l(u)+O(Kh)\).
Write \(\mathcal J\) for the tuple of root-jet bins at width
\(p\).  Given \(T,S_l(u)_p,\mathcal J\), the expansion therefore
confines the width-\(r\) bin of \((E/|E|)y_l(u)\) to at most
\(Kr^{-d}\) possibilities.  Since
\(\mathsf H(\mathcal J\mid T,S_l(u)_p)=o_*\), chain rule gives
\[
 \mathsf H\bigl(((E/|E|)y_l(u))_r\mid T,S_l(u)_p\bigr)
 \le d\log(1/r)+o_*.
\]
\Cref{a04:projection-entropy}, applied on this same event,
gives the lower bound \(\log(1/r)-o_*\).
Since \(d<1\) and \(h\) is a fixed positive power, this is impossible.
We retain dense tuple mass with
\(|F(t')-P_F(t')|\le Kh^{M+1/4}\).

\paragraph{Recovering labeled time fibers.}
Let \(z\) consist of the world, the trajectory, and the rest of the
time tree, excluding \(t'\).
The polynomial \(P_F\) is fixed at \(z\), and before restrictions
the conditional law of \(t'\) is uniform in its \(h\)-block.
If the good tuple event has mass \(p_g\ge K^{-1}\), write
\(g(z)\) for the fraction of successful \(t'\)'s in that block.
Then \(\int g(z)\,dP(z)=p_g\).
For any positive \(\delta\),
\[
 \int_{\{g<p_g\delta\}}g(z)\,dP(z)\le p_g\delta.
\]
Choose \(\delta\to0\) sufficiently slowly and retain only the other
fibers.  They still carry dense tuple mass, and every retained fiber
has successful labeled length at least \(hp_g\delta\ge h/K\).

The pre-restriction marginal law of \((l,t')\) is the original prior
times uniform time: uniform \(u\) selects \(h\)-blocks with their length
weights, and \(t'\) is uniform inside the chosen block.
Projecting the good event therefore gives a dense set of incidences
\((l,t')\).  Each has a witnessing polynomial fitting on a successful
fiber of length \(h/K\), including that incidence.
No independence is asserted after conditioning on good success or
on a discrepancy test.
If desired, round the bounded polynomial coefficients on a sufficiently
fine finite grid, adding error much smaller than \(h^{M+1/4}\), and test
its successful length on the original labels.
This makes the existential qualification test measurable and preserves
the asserted bounds.

Every construction began with upper inputs and dense mass, so it may
be performed anew on a dense residual.  This proves the proposition.
\end{proof}
